\chapter{主题索引与学习路径}\label{app:index}
\begin{longtable}{p{37mm}p{18mm}p{66mm}}\toprule 主题&位置&关联方法\\\midrule\endhead
函数与极限&\ref{sec:1:1}节&定义域、单射双射、极限定义、四则\\
无穷小与重要极限&\ref{sec:1:2}节&等价替换、误差控制\\
Taylor与洛必达&\ref{sec:1:3}节&相消阶、余项\\
Stolz与递推&\ref{sec:1:4}节&单调有界、差分\\
连续性&\ref{sec:1:5}节&零点、最值、一致连续\\
复合与隐式求导&\ref{sec:2:1}节&链式、参数、对数求导\\
高阶导数&\ref{sec:2:2}节&Leibniz\\
中值定理&\ref{sec:2:3}节&Rolle、Lagrange、Cauchy\\
曲率与极值&\ref{sec:2:4}节&凹凸、拐点\\
辅助函数&\ref{sec:2:5}节&插值、二阶估计\\
换元与双元&\ref{sec:3:1}节&根式、共轭\\
分部与部分分式&\ref{sec:3:2}节&循环积分\\
三角代换&\ref{sec:3:3}节&万能代换\\
定积分与反常积分&\ref{sec:3:4}节&区间再现、几何\\
含参与积分不等式&\ref{sec:3:5}节&Wallis、Cauchy--Schwarz\\
一阶方程&\ref{sec:4:1}节&分离、齐次、唯一性\\
线性与全微分&\ref{sec:4:2}节&伯努利、积分因子\\
降阶与齐次线性&\ref{sec:4:3}节&特征根、能量\\
非齐次方程&\ref{sec:4:4}节&共振、常数变易、算子\\
Euler与级数解&\ref{sec:4:5}节&积分方程\\
向量内积&\ref{sec:5:1}节&投影、夹角\\
叉积与混合积&\ref{sec:5:2}节&面积、体积\\
直线与平面&\ref{sec:5:3}节&方向与法向\\
空间距离&\ref{sec:5:4}节&垂足、公垂线\\
二次曲面&\ref{sec:5:5}节&截痕、切平面\\
多元可微&\ref{sec:6:1}节&路径、偏导、余项\\
多元链式隐函数&\ref{sec:6:2}节&Jacobian\\
梯度&\ref{sec:6:3}节&方向导数、切空间\\
Hessian&\ref{sec:6:4}节&驻点、鞍点\\
条件极值&\ref{sec:6:5}节&乘子、边界\\
累次积分&\ref{sec:7:1}节&Fubini、区域\\
极坐标与换序&\ref{sec:7:2}节&圆域、无初等内积分\\
三重积分&\ref{sec:7:3}节&柱面、球面坐标\\
一般换元&\ref{sec:7:4}节&Jacobian、质心\\
对称与Gaussian&\ref{sec:7:5}节&非负换序\\
曲线积分&\ref{sec:8:1}节&弧长、做功\\
Green与路径无关&\ref{sec:8:2}节&势、孔洞\\
曲面积分&\ref{sec:8:3}节&面积、通量\\
Gauss&\ref{sec:8:4}节&散度、补面\\
Stokes&\ref{sec:8:5}节&旋度、右手规则\\
正项级数&\ref{sec:9:1}节&比较、比值、根值、积分\\
任意项级数&\ref{sec:9:2}节&绝对、条件、Dirichlet、Abel\\
幂级数&\ref{sec:9:3}节&半径、端点、求和\\
一致收敛&\ref{sec:9:4}节&M判别、交换运算\\
Fourier级数&\ref{sec:9:5}节&系数、点收敛、选点求和\\
\bottomrule\end{longtable}
\section*{新增基础与应用入口}
乘法求导全表与高阶公式：第\ref{r2:product}节；含内部参数的变限求导：第\ref{r2:leibniz}节；边际与弹性：第\ref{r2:economics}节；差分方程：第\ref{r2:difference}节；积分物理建模：第\ref{r2:physics}节；投影区域：第\ref{r2:projection}节；挖洞与闭路变形：第\ref{r2:green}节；Fourier应试例题：第\ref{r2:fourier}节。
\section*{横向复习路径}
局部近似：第1章Taylor极限、第2章Taylor与曲率、第6章多元Taylor。边界与内部：第3章微积分基本定理、第8章Green、Gauss与Stokes。线性结构：第4章解空间、第5章投影、第9章函数展开。交换运算：第3章含参求导、第7章换序、第9章一致收敛。
\section*{综合自查}
能否为每个定理写出一个缺失条件的反例？能否把一道积分题转换为几何或对称问题？能否区分同一结论的局部与全局版本？能否在开始计算前说明定义域、边界、方向和收敛性？这些检查比单纯记忆公式更能减少综合题失分。
